\section{Discussion and Limitations}
\label{sec:discussion}

The multi-model evaluation and the ablation study qualify the defense-in-depth claim of the conference version in two ways. First, the effectiveness of the architecture is not uniform across models: the deterministic components, Layer~2 and Layer~3, behave identically for every model, but the prompt-level hardening, which supplies most of the block rate, depends on the instruction-following behavior of the LLM and ranges from negligible (Llama~3.2 1B) to almost complete (Nova Lite). Changing the underlying model therefore changes how early attacks are stopped, but not the leakage, which the output-side layer kept at zero for all nine models. Second, the layers are complementary in function rather than redundant in coverage: Layer~2 stops attacks before inference and never blocked a legitimate query, Layer~3 removes the residual leakage that the model and Layer~2 let through, and the false positives come from the defensive suffix, which several small models over-apply to short legitimate queries. These results argue for treating the defensive suffix as a tunable, model-specific component, and for keeping Layer~3 in every deployment regardless of the model.

From a systems perspective, the trust model exchanges hardware-backed guarantees for broader deployability. Its root of trust is a correctly configured cloud IAM, weaker than hardware attestation but available on every provider, and the defense-in-depth design mitigates this weakness: when one layer is compromised, for example through an IAM misconfiguration, the remaining layers continue to enforce protection. The architecture and TEE-based approaches are therefore complementary; recent measurements show TEE overheads below 10\% on CPUs and below 7\% on GPUs~\cite{teeperf2025}, so high-assurance deployments may combine both.

Several limitations remain. Rotating the HMAC signing key requires rebuilding the container image; risk-scoring thresholds may need tuning per deployment; text-processing pipelines that normalize Unicode may strip the zero-width watermark; adversarial fine-tuning on captured prompt--response pairs was not evaluated; and covert signaling mechanisms carry risks of their own~\cite{stegocollusion2024}. The adversarial evaluation is static: adaptive adversaries that refine their inputs from the observed blocking behavior were not evaluated, and runtime memory access lies outside the threat model of Section~\ref{sec:threats}. The new evaluations add limitations of their own. The leakage oracle detects the verbatim reproduction of 50-character chunks of the prompt and of canary tokens, the same criterion that Layer~3 applies, so a paraphrased disclosure of the instructions does not count as a leak and the outcomes classified as passed without leak must be read with this caveat, in particular for the small models that answered many attacks; an independent oracle applied to the recorded raw outputs is the natural next step. The classification of refusals relies on the fixed refusal sentence and on empty outputs, which undercounts models that decline in their own words and inflated the false-positive rate of one model whose runtime returned empty responses to long queries. The corpus is the 42-query corpus of the conference version, so each legitimate query moves the FPR by 4.3 points, and sampling-based decoding was not evaluated. The open-weight models were served at 4-bit or 8-bit quantization on a single GPU, and two of them without the demonstration tools registered, so absolute rates may differ under other serving conditions; the comparison across configurations is unaffected, since each model was evaluated under constant conditions.

From an ethical standpoint, the watermarking mechanism identifies accounts and sessions rather than individual users. Its use should be disclosed transparently and confined to intellectual property protection and forensic analysis, in line with standard privacy practices.
